\documentclass[11pt]{article}
\usepackage[T1]{fontenc}
\usepackage[utf8]{inputenc}
\usepackage[french]{babel}
\usepackage{lmodern}
\usepackage{geometry}
\geometry{margin=1in}
\usepackage{microtype}
\usepackage{amsmath,amssymb,mathtools,mathrsfs}
\usepackage{enumitem}
\usepackage{amsthm}
\usepackage{tikz-cd}
\DeclareMathOperator{\Reg}{Reg}
\DeclareMathOperator{\Supp}{Supp}
\DeclareMathOperator{\Div}{Div}
\DeclareMathOperator{\codim}{codim}
\providecommand{\cO}{\mathcal{O}}
\providecommand{\fm}{\mathfrak{m}}
\providecommand{\wt}[1]{\widetilde{#1}}
\providecommand{\Cons}{\operatorname{Cons}}
\providecommand{\Hh}{\mathrm{H}}
\providecommand{\uH}{\underline{\mathrm H}}
\providecommand{\uExt}{\underline{\Ext}}
\providecommand{\qedtext}{\hfill$\square$}
\providecommand{\ZZ}{\mathbb{Z}}
\providecommand{\QQ}{\mathbb{Q}}
\providecommand{\RR}{\mathbb{R}}
\providecommand{\FF}{\mathbb{F}}
\providecommand{\NN}{\mathbb{N}}

\newenvironment{statement}[1]{\par\medskip\noindent\textbf{#1.}\ }{\par\medskip}
\newenvironment{prooftext}[1][Proof]{\par\noindent\emph{#1.}\ }{\par\medskip}

\usepackage{hyperref}
\hypersetup{colorlinks=true,linkcolor=blue,citecolor=blue,urlcolor=blue}
\setlength{\parindent}{0pt}
\setlength{\parskip}{0.55em}
\let\accentH\H
\renewcommand{\H}{\mathrm{H}}
\DeclareMathOperator{\Hom}{Hom}
\DeclareMathOperator{\Ext}{Ext}
\DeclareMathOperator{\Tor}{Tor}
\DeclareMathOperator{\Spec}{Spec}
\DeclareMathOperator{\cd}{cd}
\newcommand{\R}{\mathrm{R}}
\newcommand{\D}{\mathrm{D}}
\newcommand{\Z}{\mathbb{Z}}
\newcommand{\F}{\mathbb{F}}
\newcommand{\Q}{\mathbb{Q}}
\newcommand{\E}{\mathrm{E}}
\newcommand{\ob}{\operatorname{ob}}
\newcommand{\car}{\operatorname{car}}
\newcommand{\ol}[1]{\overline{#1}}
\newcommand{\SGAA}{SGA~A}

\providecommand{\Ker}{\operatorname{Ker}}
\providecommand{\Coker}{\operatorname{Coker}}
\providecommand{\Gal}{\operatorname{Gal}}
\providecommand{\car}{\operatorname{car}}


\providecommand{\RHom}{\mathrm{RHom}}
\providecommand{\Dctf}{D_{\mathrm{ctf}}}
\providecommand{\Id}{\mathrm{Id}}
\providecommand{\pr}{\operatorname{pr}}
\providecommand{\Gr}{\Gamma}
\providecommand{\Fix}{\operatorname{Fix}}
\providecommand{\Loc}{\operatorname{Loc}}
\providecommand{\lto}{\longrightarrow}
\providecommand{\xto}[1]{\xrightarrow{#1}}
\providecommand{\lisom}{\xrightarrow{\sim}}
\providecommand{\ctf}{\mathrm{ctf}}
\providecommand{\et}{\mathrm{\acute{e}t}}
\providecommand{\Sfrac}{SGA~$4^{1/2}$}


\providecommand{\cl}{\operatorname{cl}}
\providecommand{\Tr}{\operatorname{Tr}}
\providecommand{\EP}{\operatorname{EP}}
\providecommand{\gr}{\operatorname{gr}}
\providecommand{\DF}{\operatorname{DF}}

\providecommand{\OO}{\mathcal{O}}
\providecommand{\LL}{\mathbf L}
\providecommand{\parf}{\mathrm{parf}}
\providecommand{\coh}{\mathrm{coh}}
\providecommand{\frob}{\operatorname{Frob}}
\nofiles

\providecommand{\RGamma}{\mathrm R\Gamma}
\providecommand{\Rj}{\mathrm Rj}
\providecommand{\Tr}{\operatorname{Tr}}
\providecommand{\id}{\operatorname{id}}
\providecommand{\cone}{\operatorname{cone}}
\providecommand{\lten}{\overset{\mathrm L}{\otimes}}
\providecommand{\Dctf}{D_{\mathrm{ctf}}}
\providecommand{\Ob}{\operatorname{ob}}


\providecommand{\Tr}{\operatorname{Tr}}
\providecommand{\RHom}{\mathrm{RHom}}
\providecommand{\Hom}{\operatorname{Hom}}
\providecommand{\ob}{\operatorname{ob}}
\providecommand{\lten}{\overset{\mathrm L}{\otimes}}
\providecommand{\nat}{\natural}
\providecommand{\dfn}{\mathrm{dfn}}


\providecommand{\Fl}{\operatorname{Fl}}
\providecommand{\MLAR}{\mathrm{MLAR}}
\providecommand{\AR}{\mathrm{AR}}
\providecommand{\NN}{\mathbb N}
\providecommand{\Ens}{\mathrm{Ens}}
\providecommand{\cC}{\mathcal C}
\providecommand{\cE}{\mathcal E}
\providecommand{\cF}{\mathcal F}
\providecommand{\cG}{\mathfrak G}
\providecommand{\source}{\operatorname{source}}
\providecommand{\op}{\mathrm{op}}


\providecommand{\End}{\operatorname{End}}
\providecommand{\id}{\operatorname{id}}
\providecommand{\Im}{\operatorname{Im}}
\providecommand{\Ker}{\operatorname{Ker}}
\providecommand{\Coker}{\operatorname{Coker}}
\providecommand{\gr}{\operatorname{gr}}

\renewcommand{\Im}{\operatorname{Im}}

\providecommand{\Rlim}{\varprojlim}
\providecommand{\adn}{\operatorname{adn}}
\providecommand{\adf}{\operatorname{adf}}
\providecommand{\AR}{\mathrm{AR}}
\providecommand{\MLAR}{\mathrm{MLAR}}
\providecommand{\Ecat}{\mathcal E}
\providecommand{\Dcat}{\mathcal D}
\providecommand{\Ccat}{\mathcal C}
\providecommand{\Bcat}{\mathcal B}
\providecommand{\Pcat}{\mathcal P}
\providecommand{\J}{J}
\providecommand{\gr}{\operatorname{gr}}
\providecommand{\grs}{\operatorname{grs}}
\providecommand{\Homcat}{\operatorname{Hom}}
\providecommand{\Ker}{\operatorname{Ker}}
\providecommand{\Coker}{\operatorname{Coker}}
\providecommand{\Im}{\operatorname{Im}}
\providecommand{\can}{\operatorname{can}}
\providecommand{\epi}{\operatorname{epi}}
\providecommand{\id}{\operatorname{id}}


\providecommand{\RGamma}{\mathrm R\Gamma}
\providecommand{\RHom}{\mathrm R\mathcal{H}om}
\providecommand{\Zell}{\mathbb Z_\ell}
\providecommand{\Qell}{\mathbb Q_\ell}
\providecommand{\Ql}{\mathbb Q_\ell}
\providecommand{\lten}{\overset{\mathrm L}{\otimes}}
\providecommand{\Ob}{\operatorname{ob}}
\providecommand{\Ens}{\operatorname{Ens}}
\providecommand{\Ensf}{\operatorname{Ensf}}
\providecommand{\Aut}{\operatorname{Aut}}
\providecommand{\End}{\operatorname{End}}
\providecommand{\Sp}{\operatorname{Sp}}
\providecommand{\Ab}{\operatorname{Ab}}
\providecommand{\Abc}{\operatorname{Abc}}
\providecommand{\Abf}{\operatorname{Abf}}
\providecommand{\fct}{\operatorname{fct}}
\providecommand{\Lc}{\operatorname{Lc}}
\providecommand{\GL}{\operatorname{GL}}
\providecommand{\car}{\operatorname{car}}
\providecommand{\cHom}{\mathcal{H}om}
\providecommand{\fc}{\operatorname{fc}}
\providecommand{\adc}{\operatorname{adc}}
\providecommand{\ad}{\operatorname{ad}}
\providecommand{\adn}{\operatorname{adn}}
\providecommand{\modn}{\operatorname{mod}}
\providecommand{\AR}{\operatorname{AR}}


\providecommand{\R}{\mathrm R}
\providecommand{\V}{\mathrm V}
\providecommand{\A}{\mathbb A}
\providecommand{\P}{\mathbb P}
\providecommand{\Z}{\mathbb Z}
\providecommand{\muS}{\mu_S}
\providecommand{\cupdot}{\smile}
\providecommand{\id}{\operatorname{id}}
\providecommand{\Tr}{\operatorname{Tr}}
\providecommand{\Spec}{\operatorname{Spec}}
\providecommand{\Supp}{\operatorname{Supp}}
\providecommand{\coker}{\operatorname{coker}}
\providecommand{\lten}{\overset{\mathrm L}{\otimes}}
\newenvironment{proposition}[1][]{\par\medskip\noindent\textbf{Proposition #1.}\ }{\par\medskip}
\newenvironment{corollaire}[1][]{\par\medskip\noindent\textbf{Corollaire #1.}\ }{\par\medskip}
\newenvironment{theoreme}[1][]{\par\medskip\noindent\textbf{Théorème #1.}\ }{\par\medskip}
\newenvironment{lemme}[1][]{\par\medskip\noindent\textbf{Lemme #1.}\ }{\par\medskip}

\begin{document}

\section*{Exposé VII}

\begin{center}
\textbf{COHOMOLOGIE DE QUELQUES SCHÉMAS CLASSIQUES}\\[0.5em]
\textbf{ET}\\[0.5em]
\textbf{THÉORIE COHOMOLOGIQUE DES CLASSES DE CHERN}\\[1em]
par\\[0.5em]
J.-P. Jouanolou
\end{center}

Dans tout le texte, on suppose donné un entier $\nu>0$ et on pose
\[
A=\mathbb Z/\nu\mathbb Z;
\]
les caractéristiques résiduelles de tous les schémas envisagés sont premières à l'entier $\nu$. Si $X$ est un tel schéma, on note $\mu_X$ le faisceau des racines $\nu$-ièmes de l'unité et $A^{\bullet}(X)$ l'anneau de cohomologie
\[
\bigoplus_i H^{2i}(X,\mu_X^{\otimes i})
\]
avec pour multiplication le cup-produit.

\subsection*{1. Fibres vectoriels}

Soient $S$ un schéma, $E$ un $\mathcal O_S$-module localement libre, sous-entendu de type fini. On pose $X=\V(E)$, $U=\V(E)^*$, complémentaire de la section nulle, et on note
\[
p:X\to S,
\qquad
q:U\to S
\]
les projections canoniques.

\begin{statement}{Proposition 1.1}
Soit $L$ un $A$-module.
\begin{enumerate}[label=(\roman*)]
\item
\begin{enumerate}[label=\alph*)]
\item Pour tout entier $j>0$, on a
\[
R^j p_*(p^*L)=0.
\]
\item Le morphisme d'adjonction
\[
L\longrightarrow p_*p^*(L)
\]
est un isomorphisme.
\end{enumerate}
\item On suppose $E$ de rang constant $r$.
\begin{enumerate}[label=\alph*)]
\item
\[
R_!^i p(p^*L)=0\quad\text{si }i\ne 2r.
\]
\item Le morphisme trace (SGA 4 XVIII)
\[
R_!^{2r}p(p^*L)\longrightarrow L\otimes\mu_S^{\otimes -r}
\]
est un isomorphisme.
\end{enumerate}
\end{enumerate}
\end{statement}

\emph{Preuve.} La partie (i) résulte immédiatement par localisation pour la topologie de Zariski de (SGA 4 XV 2.2). L'assertion (ii) est essentiellement contenue dans (SGA 4 XVIII).

\begin{statement}{Corollaire 1.2}
Soit $L$ un $A_S$-module. Pour tout entier $i$, le morphisme image réciproque
\[
H^i(S,L)\longrightarrow H^i(X,p^*L)
\]
est un isomorphisme.
\end{statement}

\emph{Preuve.} D'après (1.1), le morphisme d'adjonction $L\to Rp_*(p^*L)$ est un isomorphisme. Comme $H^*(X,p^*L)=H^*(S,Rp_*p^*(L))$, l'assertion en résulte.

\begin{statement}{Proposition 1.3}
Supposons $E$ de rang constant $r$. Pour tout $A_S$-module $L$, les propriétés suivantes sont vérifiées.
\begin{enumerate}[label=(\roman*)]
\item
\begin{enumerate}[label=\alph*)]
\item
\[
R^i q_*(q^*L)=0\quad\text{si }i\ne 0,2r-1.
\]
\item Le morphisme d'adjonction
\[
L\longrightarrow q_*q^*(L)
\]
est un isomorphisme.
\item Désignant par $Y$ le complémentaire de $U$ dans $X$, le morphisme
\[
R^{2r-1}q_*(q^*L)\longrightarrow L\otimes\mu_S^{\otimes -r}
\]
composé du morphisme bord
\[
R^{2r-1}q_*(q^*L)\longrightarrow R_Y^{2r}p_*(p^*L)
\tag{1}
\]
et de l'isomorphisme de pureté (SGA 4 XVIII)
\[
R_Y^{2r}p_*(p^*L)\longrightarrow L\otimes\mu_S^{\otimes -r}
\]
est un isomorphisme.
\end{enumerate}
\item
\begin{enumerate}[label=\alph*)]
\item
\[
R_!^i q(q^*L)=0\quad\text{si }i\ne 1,2r.
\]
\item Le morphisme bord
\[
L\longrightarrow R_!^1 q(q^*L)
\]
provenant de la suite exacte illimitée, où $f:Y\to S$ est la projection canonique,
\[
\cdots\longrightarrow R_!^i q(q^*L)\longrightarrow R_!^i p(p^*L)
\longrightarrow R^i f_*(f^*L)\longrightarrow R_!^{i+1}q(q^*L)\longrightarrow\cdots
\]
est un isomorphisme.
\item Le morphisme trace
\[
R_!^{2r}q(q^*L)\longrightarrow L\otimes\mu_S^{\otimes -r}
\]
est un isomorphisme.
\end{enumerate}
\end{enumerate}
\end{statement}

\emph{Preuve.} Soient $\alpha:U\to X$ et $\beta:Y\to X$ les immersions canoniques. On sait (SGA 4 XVI 3.7) que, posant $F=p^*(L)$, on a
\[
F\xrightarrow{\sim}\alpha_*\alpha^*(F)
\quad\text{et}\quad
(R^j\alpha_*)\alpha^*(F)=0\quad\text{si }j\ne0,2r-1.
\]
Par ailleurs les faisceaux $(R^j\alpha_*)\alpha^*F$ pour $j\ge1$ sont à support dans $Y$, donc acycliques pour le foncteur $p_*$. On déduit de ces remarques les assertions (i) a) et b) en utilisant la suite spectrale de Leray
\[
R^i p_*\bigl(R^j\alpha_*\bigr)(q^*L)\Longrightarrow R^n q_*(q^*L).
\]
L'assertion (i) c) résulte de (1.1 (i)) et de la suite exacte illimitée
\[
\cdots\longrightarrow R_Y^i p_*(p^*L)\longrightarrow R^i p_*(p^*L)\longrightarrow R^i q_*(q^*L)
\longrightarrow R_Y^{i+1}p_*(p^*L)\longrightarrow\cdots.
\]
Le faisceau $R_Y^{2r}p_*(p^*L)$ est associé au préfaisceau
\[
T\longmapsto H^{2r}_{Y\times_S T}(X\times_S T,p^*L).
\]

Prouvons (ii). L'assertion c) provient de ce que le morphisme $q$ est lisse et a ses fibres géométriques connexes de dimension $r$. La suite exacte mentionnée en b) est obtenue en appliquant le foncteur $R_!p$ à la suite exacte canonique
\[
0\longrightarrow \alpha_!(q^*L)\longrightarrow p^*L\longrightarrow \beta_!(\beta^*p^*L)\longrightarrow 0.
\]
Utilisant (1.1 (i)), on déduit sans peine de cette suite exacte illimitée les assertions (ii) a) et b).

\begin{statement}{Remarque 1.4}
Pour $S=\operatorname{Spec}(\mathbb C)$, (1.3) exprime que le complémentaire de l'origine dans $\mathbb C^r$ a même cohomologie que la sphère réelle de dimension $2r-1$, rétracte par déformation de $\mathbb C^r-0$.
\end{statement}

\begin{statement}{Corollaire 1.5 (suite exacte de Gysin)}
La suite spectrale de Leray
\[
H^i(S,R^j q_*(q^*L))\Rightarrow H^{i+j}(U,q^*L)
\]
définit une suite exacte illimitée, appelée suite exacte de Gysin,
\[
\cdots\longrightarrow H^i(S,L)\longrightarrow H^i(U,q^*L)\longrightarrow
H^{i-2r+1}(S,L\otimes\mu_S^{\otimes -r})\longrightarrow H^{i+1}(S,L)\longrightarrow\cdots.
\]
\end{statement}

\emph{Preuve.} Lemme des deux lignes.

\subsection*{2. Schémas projectifs}

Dans ce qui suit, le mode d'exposition adopté est basé sur une suggestion de L. Illusie, et la méthode suivie est voisine d'une méthode de P. Deligne pour prouver la dégénérescence de certaines suites spectrales [6].

\subsubsection*{2.1}
Soient $X$ et $Y$ deux schémas et $f:X\to Y$ un morphisme. Nous allons définir une flèche
\[
Rf_*(A_X)\otimes Rf_*(A_X)\longrightarrow Rf_*(A_X)
\tag{2.1.1}
\]
de $D(Y)$, ayant les propriétés d'associativité et de commutativité usuelles, et qui munit donc $Rf_*(A_X)$ d'une structure d'algèbre généralisée. Notons pour cela
\[
C^*(X,A_X)
\]
la résolution de Godement de $A_X$ (SGA 4 XVII). Procédant comme dans (TF 6.6), on définit un morphisme de résolutions
\[
C^*(X,A_X)\otimes C^*(X,A_X)\longrightarrow C^*(X,A_X\otimes A_X).
\]
En le composant avec le morphisme
\[
C^*(X,A_X\otimes A_X)\longrightarrow C^*(X,A_X)
\]
déduit de la multiplication de $A_X$, on obtient enfin un morphisme
\[
C^*(X,A_X)\otimes C^*(X,A_X)\longrightarrow C^*(X,A_X).
\tag{2.1.2}
\]
Appliquant le foncteur $f_*$, on en déduit sans peine un morphisme de complexes de $A_Y$-modules
\[
f_*C^*(X,A_X)\otimes f_*C^*(X,A_X)\longrightarrow f_*C^*(X,A_X),
\]
d'où aussitôt (2.1.1) en utilisant des résolutions plates des composants de l'expression de gauche.

\subsubsection*{2.2}
Soient $S$ un schéma et $E$ un $\mathcal O_S$-module localement libre de rang $r+1$. On note
\[
P=P_S(E)\xrightarrow{p}S
\]
le fibré projectif correspondant. Le $\mathcal O_P$-module canonique $\mathcal O_P(1)$ définit, comme on sait, grâce à la suite exacte de Kummer, un élément
\[
\xi\in H^2(P,\mu)
\]
et on notera $\eta$ l'image de $\xi$ par le morphisme canonique
\[
H^2(P,\mu)\longrightarrow H^0(S,R^2p_*(\mu))
\]
provenant de la suite spectrale de Leray
\[
H^i(S,R^jp_*(\mu))\Rightarrow H^{i+j}(P,\mu).
\]
La classe $\xi$ s'identifie à une flèche $A_P\to\mu_P[2]$ de $D(P)$, qui définit par transposition $\mu_P^{\otimes -1}[-2]\to A_P$, d'où par l'adjonction usuelle une flèche
\[
c:\mu_S^{\otimes -1}[2]\longrightarrow Rp_*(A_P),
\]
coïncidant avec $\eta$ sur la cohomologie de dimension $2$.

Par produit tensoriel, la flèche $c$ définit, grâce à (2.2.1), des flèches
\[
c_i:\mu_S^{\otimes -i}[-2i]\longrightarrow Rp_*(A_P),
\]
pour tout entier $i\ge1$, et on convient de noter $c_0$ la flèche d'adjonction
\[
A_S\longrightarrow Rp_*(A_P).
\]
On définit ainsi une flèche
\[
\gamma=\bigoplus_{i=0}^{r}c_i:
\bigoplus_{i=0}^{r}\mu_S^{\otimes -i}[-2i]\longrightarrow Rp_*(A_P)
\]
de $D(S)$.

\begin{statement}{Théorème 2.2.1}
La flèche $\gamma$ est un isomorphisme.
\end{statement}

Une fois définie la flèche $\gamma$, (2.2.1) équivaut à son corollaire suivant.

\begin{statement}{Proposition 2.2.2}
Sous les hypothèses précédentes, on a:
\begin{enumerate}[label=\alph*)]
\item
\[
R^{2i+1}p_*(A_P)=0\quad\text{pour tout }i\in\mathbb N.
\]
\item Lorsqu'on munit $\bigoplus_i R^{2i}p_*(\mu_P^{\otimes i})$ de sa structure de $A_S$-algèbre déduite du cup-produit, le morphisme de $A_S$-algèbres
\[
A_S[T]\longrightarrow \bigoplus_i R^{2i}p_*(\mu_P^{\otimes i})
\]
défini par $\eta$ a pour noyau l'idéal $(T^{r+1})$ et définit donc un isomorphisme de $A_S$-algèbres
\[
\theta:A_S[T]/(T^{r+1})\xrightarrow{\sim}\bigoplus_i R^{2i}p_*(\mu_P^{\otimes i}).
\]
\item Le morphisme
\[
A_S\longrightarrow R^{2r}p_*(\mu^{\otimes r})
\]
obtenu en composant $\theta$ et la restriction à $A_S$ de la multiplication par $T^r$ de $A_S[T]/(T^{r+1})$ dans lui-même est l'inverse du morphisme trace (SGA 4 XVIII)
\[
R^{2r}p_*(\mu^{\otimes r})\longrightarrow A_S;
\]
en d'autres termes, on a
\[
\operatorname{Tr}(\eta^r)=1.
\]
\end{enumerate}
\end{statement}

\emph{Preuve.} Les données étant stables par changement de base, le théorème de changement de base propre (SGA 4 XII 5.1), et la compatibilité du morphisme trace avec changement de base (SGA 4 XVIII), permettent de se ramener à vérifier l'assertion analogue pour les fibres de $P$. Autrement dit, on peut supposer que $S$ est le spectre d'un corps séparablement clos $k$ et que $p$ est le morphisme canonique $\mathbb P^r_k\to k$. Dans ce cas, la démonstration utilise le lemme suivant.

\begin{statement}{Lemme 2.2.3}
Soient $k$ un corps séparablement clos et $p:\mathbb P_k^r\to k$ la projection canonique. Alors:
\begin{enumerate}[label=\alph*)]
\item
\[
R^{2i+1}p_*(A_P)=0\quad\text{pour tout }i.
\]
\item
\[
R^{2i}p_*(A_P)\simeq A_k\quad\text{pour }i\in[0,r],
\qquad
R^{2i}p_*(A_P)=0\quad\text{si }i>r.
\]
\end{enumerate}
\end{statement}

Le lemme, manifestement vrai pour $r=0$, se montre par récurrence sur $r$. Supposons donc qu'il soit vrai pour $r-1$ et montrons-le pour $r$. L'espace projectif $\mathbb P_k^{r-1}$ se plonge comme fermé dans $\mathbb P_k^r$ avec pour complémentaire l'ouvert $\mathbb E_k^r$, d'où une suite exacte de cohomologie à supports propres
\[
\cdots\longrightarrow H_!^i(\mathbb E_k^r)\longrightarrow H^i(\mathbb P_k^r)
\longrightarrow H^i(\mathbb P_k^{r-1})\longrightarrow H_!^{i+1}(\mathbb E_k^r)\longrightarrow\cdots.
\]
Le calcul de la cohomologie à supports propres de $\mathbb E_k^r$ (1.1) montre que, si $i\ne(2r,2r-1)$, l'image réciproque
\[
H^i(\mathbb P_k^r)\longrightarrow H^i(\mathbb P_k^{r-1})
\]
est un isomorphisme, d'où l'assertion dans ce cas. Si $i=2r-1$, on voit de même que $H^{2r-1}(\mathbb P_k^r)$ est contenu dans $H^{2r-1}(\mathbb P_k^{r-1})$, donc nul par hypothèse de récurrence. Enfin, si $i=2r$, les relations
\[
H^{2r-1}(\mathbb P_k^{r-1})=H^{2r}(\mathbb P_k^{r-1})=0
\]
montrent que
\[
H_!^{2r}(\mathbb E_k^r)\xrightarrow{\sim}H^{2r}(\mathbb P_k^r),
\]
d'où l'assertion d'après (1.1).

Le lemme entraîne la partie a) de 2.2.1. Prouvons maintenant que $\theta$ est un isomorphisme. Le morphisme $\theta$ peut être interprété comme un morphisme d'anneaux
\[
A[T]/(T^{r+1})\longrightarrow\bigoplus_i H^{2i}(P,\mu^{\otimes i}).
\]

\end{document}
